% Printed 441.
\par\clearpage
\noindent\null\hfill{\small\textbf{HEIGHT MEASUREMENTS.}}\hfill\makebox[0pt][r]{\fontsize{13}{16}\selectfont 441}\par\medskip
\par\noindent\hspace*{\parindent}The whole system of this year's measurements comprises 29 lines. Two of these, however,
already belong to the work of 1822 and 1823 as far as the measurements of the reciprocal
zenith distances are concerned (namely Falkenberg-Wilsede and Wilsede-Hamburg).
On two lines the zenith distance was measured in one direction \textls[120]{only} (there was no
measurement at all from Bremen to Bottel or from Brüttendorf to Zeven). On two others
it was also measured in one direction only in the sense that, for lack of time, the
theodolite alone was used for the measurement from the other end (namely from Bottel
to Steinberg and from Wilsede to Bullerberg\textsuperscript{*}\,). Finally, on the Litberg-Hamburg
line\textsuperscript{**} the target in Hamburg was the ball, whose height above the windows of
the chamber I do not yet know exactly. This leaves 22 lines on which the reciprocal
zenith distances were fully measured, i.e. with the 12-inch circle and with at least
20 repetitions, and on which heliotrope light was the target \textls[120]{without exception}.
The sum of the 22 curvatures (amplitudes) is \(6^\circ15'49^{\prime\prime}844\) (the largest,
from Wilsede to Steinberg, is \(26'4^{\prime\prime}{,}883\)), while the sum of the 22 refractions
is \(55'32^{\prime\prime}{,}329\). Thus the mean ratio is as 1 to
\[0{,}14778.\]
\par\noindent\hspace*{\parindent}The mean of the 22 individual quotients is \(0{,}14499\). I prefer the former result
because the quotient fluctuates much more for \textls[120]{small} arcs than for larger ones.
I have already noted elsewhere that by refraction I mean the difference between the
directions at the two endpoints. Most writers, on the other hand, call half that
difference refraction, namely the angles that the tangents to the light ray's path
at the two endpoints make with the chord. Under that definition, then, this year's
measurements give \(0{,}07389\). All my measurements from 1821---1823 that were usable for
this purpose had given, for 34 lines, a sum of the curvatures of \(10^\circ8'14^{\prime\prime}{,}249\),
a sum of the refractions \(=1^\circ17'43^{\prime\prime}{,}670\), and thus a ratio as 1 to \(0{,}12779\).
The longest line among them is Lichtenberg-Falkenberg, whose curvature is
\(46'7^{\prime\prime}{,}250\). One would be quite mistaken in supposing that the larger
quotient of 1824 was peculiar to the flatter districts. In fact the measurements
of 1823, which belong to equally flat districts, give an even smaller result than
all those of
\par\medskip\noindent\hspace*{0.03\linewidth}{\footnotesize IX.}\hfill\makebox[0.18\linewidth][l]{56}\par

% Printed 442.
\par\clearpage
\noindent\makebox[0pt][l]{\fontsize{13}{16}\selectfont 442}\hfill{\small\textbf{CORRESPONDENCE.}}\hfill\null\par\medskip
\noindent 1821---1823. I have, however, no doubt at all about the source of the difference.
In the years 1821---1823 I measured almost all the zenith distances in the morning,
at noon, or soon after noon, since I considered those measurements a secondary matter.
As a rule I devoted to them only the hours when visibility was insufficient for the
theodolite measurements. On \textls[120]{sunny days} in flat country these are invariably
the hours just mentioned, and on sunny days the refraction is then invariably smaller
than somewhat later in the afternoon, when the images become steadier and sharper.
In 1824, however, the great majority of my lines ran so close above intervening
obstacles that in the morning no light at all could get through on very many of
them. With few exceptions I therefore gave up morning measurements of the principal
directions altogether. The zenith distances in 1824 were therefore \textls[120]{all}
measured in the afternoon in already more favourable air, though never shortly before
sunset, when in all my experience the refraction becomes considerably greater still.
They were mostly measured in the afternoon from 3 to 4 o'clock\textsuperscript{*}), and I
tried, as far as practicable, to measure the reciprocal zenith distances on each
line under approximately the same atmospheric conditions.
\par\noindent\hspace*{\parindent}If you would like to have all my \(22+34=56\) results individually, they are gladly
at your disposal. For the excluded lines marked \textsuperscript{* **} above, the values are
\begin{center}
\setlength{\tabcolsep}{4pt}
\begin{tabular}{lrrr}
& Curvatures & Refractions & Quotient\\
{}[Wilsede-Bullerberg]\textsuperscript{*} &\(18'56^{\prime\prime}{,}002\)&\(2'42^{\prime\prime}{,}975\)&\(0{,}14346\)\\
{}[Litberg-Hamburg]\textsuperscript{**} &\(16\phantom{' }14{,}047\)&\(2\phantom{' }49{,}295\)&\(0{,}17380\)\\\cline{2-4}
\multicolumn{1}{r}{Total}&\(35'10^{\prime\prime}{,}049\)&\(5'32^{\prime\prime}{,}270\)&\(0{,}15747\),
\end{tabular}
\end{center}
\noindent taking, for \textsuperscript{**}, the height of the ball in Hamburg above the windows
to be \(=15{,}558\,\mathrm{m}\). You would oblige me by determining these relative heights,
and the height of your barometer in Altona relative to the two points above,
more accurately.\quad\mbox{.\ .\ .\ .\ .\ .}
\par\medskip\noindent\rule{.25\linewidth}{.4pt}\par
{\small\noindent\hspace*{\parindent}*) As long as the days were still longer. Later, gradually somewhat earlier.
In October the air usually had little shimmer left by \(1\frac{1}{2}\) or 2 o'clock.\par}
